\begin{proposition}[6.13.4]
\label{IV.2.s6.13.4-fr}
Sea $X$ un preesquema localmente noetheriano. Para que
$\operatorname{Nor}(X)$ sea abierto, es necesario y suficiente que todo
punto maximal $x$ de $X$ tal que $x\in\operatorname{Nor}(X)$ \textup{(lo que
significa que $X$ es \emph{reducido} en el punto $x$)} sea un punto interior de
$\operatorname{Nor}(X)$.
\end{proposition}

\begin{proof}
Solo hay que demostrar la suficiencia de la condición. En virtud de la
observación de (6.13.1), es posible limitarse (reemplazando $X$ por un
abierto de $X$ en el que $X$ es reducido) al caso en que $X$ es \emph{reducido},
es decir, suponer que los puntos maximales de $X$ pertenecen a
$U_{S_2}(X)$ y a $U_{R_1}(X)$; como
$\operatorname{Nor}(X)=U_{S_2}(X)\cap U_{R_1}(X)$ por el criterio de
Serre (5.8.6), se deduce de (6.11.7) y (6.12.10) que ambos conjuntos
son abiertos; por consiguiente, $\operatorname{Nor}(X)$ es abierto en $X$.
\end{proof}
